% GENERATED FILE. Edit canonical lessons, then run npm run generate:books.
% canonical-source-hash: fnv1a64:b9bf2d170686fe7f
% canonical-lessons: FR-C24-prendre, FR-C24-demander, FR-C24-aider, FR-C24-aimer

\chapter{Taking, Asking, Helping, Loving}
\label{ch:fr-taking-asking-helping-loving}
\hlchaptermodality{24}

\begin{chapteropening}
\textbf{By the end of this chapter:} \emph{I can say that I take, ask, help and like or love in French, break prendre's stem the way its plural requires, and know that je t'aime bien is weaker than je t'aime rather than stronger.}

Run the chapter's four verbs together with je — je prends, je demande, j'aide, j'aime — then turn aimer over: one verb covers like and love, the thing named decides which, and adding bien downgrades it. Reaches back past the chapter to re-say le chien and le chat from Chapter 22, le vert from Chapter 23, and the s'il vous plaît request from Chapter 19.
\end{chapteropening}

\section[prendre]{prendre — \textquotedblleft{}to take,\textquotedblright{} the verb your hand does}

\label{lesson:FR-C24-prendre}



\begin{quote}
You already have \textbf{la main}, the hand. Here is the verb that hand performs — and it turns out to be hiding inside a surprising number of English words you use without thinking.
\end{quote}

\subsection*{You'll want to know: the word}
\begin{quote}
\textbf{prendre} — to take. Said \emph{PRAHN-druh}: the \emph{-en-} is the \textbf{nasal} vowel of \emph{an}, sent through the nose, not a clean \emph{en}.
\end{quote}

>

\begin{quote}
\textbf{Je prends le pain.} — I take the bread.
\end{quote}

The \emph{-ds} of \emph{je prends} is \textbf{silent}. You say \emph{prahn}, and stop.

\begin{cousinweb}
\textbf{prendre} $\leftarrow$ Latin \textbf{prehendere}, \textquotedblleft{}to grasp, to seize,\textquotedblright{} which Romans themselves shortened to \textbf{prendere}. That short form is the one French inherited — and English borrowed the family repeatedly:

\begin{itemize}
\raggedright
  \item \textbf{prehensile} — a tail that grasps.
  \item \textbf{apprehend} — to seize (\emph{ad-} \textquotedblleft{}toward\textquotedblright{} + grasp); \textbf{reprehensible}, something to be seized back on.
  \item \textbf{prison} $\leftarrow$ Latin \emph{prehensiō}, \textquotedblleft{}\textbf{a seizing}\textquotedblright{} — a jail is, literally, a taking-hold.
  \item \textbf{surprise} — French \emph{sur-} \textquotedblleft{}over\textquotedblright{} + \emph{prendre}: an \textbf{over-taking}.
  \item \textbf{enterprise}, something undertaken; \textbf{apprentice}, one who takes on.
  \item \textbf{impregnable} — from Old French \emph{imprenable}, \textquotedblleft{}\textbf{not takeable}.\textquotedblright{} Nothing to do with pregnancy.
\end{itemize}

\emph{Manus} gave you a hand that \textbf{holds} — \emph{manual}, \emph{maintain}. \emph{Prendere} gives you the hand that \textbf{closes}.
\end{cousinweb}

\begin{grammarlens}[title={Grammar Lens: the stem that doubles}]
\emph{Prendre} is irregular, and its irregularity is one clean rule: the stem \textbf{loses} its \emph{d}, and in the \emph{ils} form it \textbf{doubles its n}. Say the six:

\begin{itemize}
\raggedright
  \item \textbf{je prends} — \emph{prahn}
  \item \textbf{tu prends} — \emph{prahn}
  \item \textbf{il prend} — \emph{prahn} (all three identical to the ear)
  \item \textbf{nous prenons} — \emph{pruh-NOHN}, and the nasal is \textbf{gone}
  \item \textbf{vous prenez} — \emph{pruh-NAY}, nasal gone again
  \item \textbf{ils prennent} — \emph{PREN}, a doubled \emph{n} and a flat \emph{e}
\end{itemize}

So the nasal survives in the singular and dies in the plural. That is the whole verb.
\end{grammarlens}

\subsection*{Your turn}
Say these aloud:

\begin{itemize}
\raggedright
  \item \textquotedblleft{}prendre\textquotedblright{} — \emph{PRAHN-druh}, nasal in the middle
  \item \textquotedblleft{}je prends le pain\textquotedblright{} then \textquotedblleft{}nous prenons le pain\textquotedblright{} — nasal, then no nasal
  \item \textquotedblleft{}je prends la main\textquotedblright{} — I take the hand, the \emph{manus} you already have
  \item the family — \textquotedblleft{}prehensile, apprehend, prison, surprise\textquotedblright{} — all a \textbf{grasp}
\end{itemize}

\begin{tcolorbox}[breakable,colback=teal!4,colframe=teal!35!black,title={Before you move on}]
Say \textquotedblleft{}I take the bread.\textquotedblright{} (\emph{Je prends le pain} — and the \emph{-ds} is silent.) What happens to the nasal in \emph{nous prenons}? (It \textbf{disappears}; the stem flattens to \emph{pruh-}.) What did Latin \emph{prehendere} mean, and what is a \emph{prison} literally? (\textbf{To grasp}; a prison is \textquotedblleft{}\textbf{a seizing}\textquotedblright{}.) And the hand that does the taking? (\emph{La main}.)
\end{tcolorbox}
\section[demander]{demander — \textquotedblleft{}to ask,\textquotedblright{} and the trap it sets in English}

\label{lesson:FR-C24-demander}



\begin{quote}
\emph{Prendre} was taking. This one is the opposite move: putting a question out. And it is the word most likely to make you sound furious when you meant to sound polite.
\end{quote}

\subsection*{You'll want to know: the word}
\begin{quote}
\textbf{demander} — to ask. Said \emph{duh-mahn-DAY}: the \emph{-an-} is that same nasal from \emph{prendre}, and the \emph{-er} ending is a clean \emph{ay}.
\end{quote}

>

\begin{quote}
\textbf{Je demande l'heure.} — I ask the time.
\end{quote}

It is an ordinary \textbf{-er} verb, so nothing in it breaks: \emph{je demande, tu demandes, il demande, nous demandons, vous demandez, ils demandent}. Only \emph{nous} and \emph{vous} are audible; the other four all sound \emph{duh-MAHND}.

\begin{cousinweb}
\textbf{demander} $\leftarrow$ Latin \textbf{dēmandāre}, \textquotedblleft{}to hand over, to entrust\textquotedblright{} — \emph{dē-} plus \textbf{mandāre}. And \emph{mandāre} is itself two words: \emph{\textbf{manus}}, \textquotedblleft{}\textbf{hand},\textquotedblright{} plus \emph{\textbf{dare}}, \textquotedblleft{}\textbf{to give}.\textquotedblright{} Giving into the hand.

So the \emph{manus} you met in \emph{la main} is sitting inside this verb too:

\begin{itemize}
\raggedright
  \item \textbf{mandate} — what was placed in someone's hands to carry out.
  \item \textbf{command} — \emph{com-} + \emph{mandāre}: to entrust an order; and \textbf{commandment}, ten of them.
  \item \textbf{mandatory} — handed over as an obligation; \textbf{remand}, handed back.
\end{itemize}

One family, two very different jobs: \emph{manual} is work done \textbf{by} the hand, \emph{mandate} is a charge given \textbf{into} it.
\end{cousinweb}

\begin{culture}
\begin{quote}
\textbf{Careful.} French \textbf{demander} means simply \textquotedblleft{}\textbf{to ask}.\textquotedblright{} It is neutral, everyday and perfectly polite. English \textbf{demand} is forceful — a thing you insist on.
\end{quote}

\emph{Je demande l'heure} is \textquotedblleft{}I'm asking the time,\textquotedblright{} not \textquotedblleft{}I demand the time.\textquotedblright{} The English word drifted toward force; the French one never did. Translate \emph{demander} as \textbf{ask}, every time, and you will never be rude by accident.
\end{culture}

\subsection*{Your turn}
Say these aloud:

\begin{itemize}
\raggedright
  \item \textquotedblleft{}demander\textquotedblright{} — \emph{duh-mahn-DAY}, nasal in the middle
  \item \textquotedblleft{}je demande l'heure\textquotedblright{} — I \textbf{ask} the time, not demand it
  \item \textquotedblleft{}je prends, je demande\textquotedblright{} — take, then ask
  \item the family — \textquotedblleft{}mandate, command, mandatory\textquotedblright{} — a \textbf{hand} and a \textbf{giving}
\end{itemize}

\begin{tcolorbox}[breakable,colback=teal!4,colframe=teal!35!black,title={Before you move on}]
Say \textquotedblleft{}I ask the time.\textquotedblright{} (\emph{Je demande l'heure}.) Does \emph{demander} carry the force of English \emph{demand}? (\textbf{No} — it is the ordinary, polite word for asking.) What two Latin words are inside \emph{mandāre}, and which one do you already know? (\emph{\textbf{Manus}} \textquotedblleft{}hand\textquotedblright{} — the source of \emph{la main} — plus \emph{\textbf{dare}}, \textquotedblleft{}to give.\textquotedblright{})
\end{tcolorbox}
\section[aider]{aider — \textquotedblleft{}to help,\textquotedblright{} and the English word that is really the same word}

\label{lesson:FR-C24-aider}



\begin{quote}
You can ask for something now. This is the verb for what you hope the other person then does — and it is one of the rare cases where the English cousin is not a distant relative but the very same word, spelled shorter.
\end{quote}

\subsection*{You'll want to know: the word}
\begin{quote}
\textbf{aider} — to help. Said \emph{ay-DAY}: the \emph{ai-} is a plain \emph{ay}, the \emph{-er} ending is another \emph{ay}. Two of them in a row.
\end{quote}

>

\begin{quote}
\textbf{Vous m'aidez.} — You are helping me.
\end{quote}

Another ordinary \textbf{-er} verb: \emph{j'aide, tu aides, il aide, nous aidons, vous aidez, ils aident}. The \emph{j'} is the elision you already make in \emph{j'ai}.

\begin{cousinweb}
\textbf{aider} $\leftarrow$ Latin \textbf{adiūtāre}, a repeated-action form of \textbf{adiuvāre} — \emph{ad-} \textquotedblleft{}toward\textquotedblright{} plus \textbf{iuvāre}, \textquotedblleft{}\textbf{to help}.\textquotedblright{} English took the same verb through French and clipped it:

\begin{itemize}
\raggedright
  \item \textbf{aid} — the noun and verb; \textbf{aide}, a person who is your help.
  \item \textbf{aide-de-camp} — literally a \textquotedblleft{}camp helper,\textquotedblright{} an officer's assistant.
  \item \textbf{adjutant}, the officer who assists; \textbf{adjuvant}, a substance in medicine that helps another one work.
\end{itemize}

Two honesty notes. English \textbf{help} is \textbf{Germanic} and has no connection to \emph{aider} whatsoever. And \textbf{juvenile} is \emph{not} in this family: that is Latin \emph{iuvenis}, \textquotedblleft{}young,\textquotedblright{} a different word that merely resembles \emph{iuvāre}.
\end{cousinweb}

\begin{grammarlens}[title={Grammar Lens: asking for it politely}]
Put the verb together with the courtesy you already have:

\begin{quote}
\textbf{Vous m'aidez, s'il vous plaît.} — Help me, please. \textbf{Tu m'aides, s'il te plaît.} — the same request, to a friend.
\end{quote}

The register has to match on \textbf{both} halves. \emph{Vous} takes \emph{s'il vous plaît}; \emph{tu} takes \emph{s'il te plaît}. Mixing them — \emph{tu m'aides, s'il vous plaît} — is the mistake a beginner makes most, and a French ear catches it instantly.
\end{grammarlens}

\subsection*{Your turn}
Say these aloud:

\begin{itemize}
\raggedright
  \item \textquotedblleft{}aider\textquotedblright{} — \emph{ay-DAY}, two \emph{ay} sounds
  \item \textquotedblleft{}j'aide\textquotedblright{} then \textquotedblleft{}vous aidez\textquotedblright{}
  \item \textquotedblleft{}vous m'aidez, s'il vous plaît\textquotedblright{} — both halves formal
  \item \textquotedblleft{}tu m'aides, s'il te plaît\textquotedblright{} — both halves familiar
  \item \textquotedblleft{}je demande, vous aidez\textquotedblright{} — I ask, you help
\end{itemize}

\begin{tcolorbox}[breakable,colback=teal!4,colframe=teal!35!black,title={Before you move on}]
Ask a stranger for help, politely. (\emph{Vous m'aidez, s'il vous plaît.}) And a friend? (\emph{Tu m'aides, s'il te plaît} — both halves change together.) What Latin verb is \emph{aider} from, and which English word is the same word? (\textbf{Adiūtāre}; English \textbf{aid}.) Is English \emph{help} related? (\textbf{No} — it is Germanic, unrelated.)
\end{tcolorbox}
\section[aimer]{aimer — one verb for \textquotedblleft{}like\textquotedblright{} and \textquotedblleft{}love,\textquotedblright{} and the word that weakens it}

\label{lesson:FR-C24-aimer}



\begin{quote}
Three verbs so far: \emph{je prends}, \emph{je demande}, \emph{j'aide}. The fourth does a job English needs \textbf{two} words for — and the way you soften it is the opposite of what you would guess.
\end{quote}

\subsection*{You'll want to know: the word}
\begin{quote}
\textbf{aimer} — to like \textbf{and} to love. Said \emph{ay-MAY}: the \emph{ai-} is \emph{ay}, the \emph{-mer} is \emph{may}.
\end{quote}

>

\begin{quote}
\textbf{J'aime le chien.} — I like the dog. \textbf{J'aime le chat.} — I like the cat.
\end{quote}

An ordinary \textbf{-er} verb: \emph{j'aime, tu aimes, il aime, nous aimons, vous aimez, ils aiment}. French does not choose between \textquotedblleft{}like\textquotedblright{} and \textquotedblleft{}love\textquotedblright{} — the \textbf{thing you name} does. A dog, a colour, a food: that reads as \emph{like}. A person: \emph{love}.

\begin{cousinweb}
\textbf{aimer} $\leftarrow$ Latin \textbf{amāre}, \textquotedblleft{}to love,\textquotedblright{} and English is full of it:

\begin{itemize}
\raggedright
  \item \textbf{amateur} — one who does a thing \textbf{for love}, not for money.
  \item \textbf{amorous}, full of love; \textbf{enamoured}, in it; \textbf{paramour}, beside it.
  \item \textbf{amiable} and \textbf{amicable} — from \emph{amīcus}, \textquotedblleft{}friend\textquotedblright{}; \textbf{amity} between nations.
\end{itemize}

An honest stop, though: unlike \emph{viridis} behind \textbf{vert}, \emph{amāre} has \textbf{no agreed ancestor} beyond Latin. It is often guessed to be a nursery word, a baby's \emph{amma}, but that is a guess. The trail genuinely ends at Latin.
\end{cousinweb}

\begin{culture}
\begin{quote}
\textbf{Je t'aime.} — I love you. Full strength. \textbf{Je t'aime bien.} — I like you. \textbf{Much} weaker.
\end{quote}

Every English speaker gets this backwards. \emph{Bien} means \textquotedblleft{}well,\textquotedblright{} so \emph{je t'aime bien} looks like it ought to be \textbf{more}. It is \textbf{less}. Adding \emph{bien} — or \emph{beaucoup} — downgrades love to friendly liking, which is exactly why \emph{je t'aime bien} is the polite way to turn someone down.

With a thing, no such danger: \emph{j'aime bien le chat} is simply \textquotedblleft{}I quite like the cat.\textquotedblright{} The trap only bites when the object is a \textbf{person}.
\end{culture}

\subsection*{Your turn}
Say these aloud:

\begin{itemize}
\raggedright
  \item \textquotedblleft{}j'aime le chien\textquotedblright{} then \textquotedblleft{}j'aime le chat\textquotedblright{} — liking, because they are animals
  \item \textquotedblleft{}j'aime le vert\textquotedblright{} — I like green
  \item \textquotedblleft{}je t'aime\textquotedblright{} then \textquotedblleft{}je t'aime bien\textquotedblright{} — love, then only liking
  \item the four verbs — \textquotedblleft{}je prends, je demande, j'aide, j'aime\textquotedblright{}
  \item \textquotedblleft{}vous m'aidez, s'il vous plaît\textquotedblright{} — the request, once more
  \item \textquotedblleft{}aimer, amateur, amiable\textquotedblright{} — the \emph{amāre} family, which stops at Latin
\end{itemize}

\begin{tcolorbox}[breakable,colback=teal!4,colframe=teal!35!black,title={Before you move on}]
Which is stronger, \emph{je t'aime} or \emph{je t'aime bien}? (\emph{\textbf{Je t'aime}} — \emph{bien} makes it \textbf{weaker}.) Say \textquotedblleft{}I like the dog\textquotedblright{} and \textquotedblleft{}I like green.\textquotedblright{} (\emph{J'aime le chien. J'aime le vert}.) What does \emph{amateur} literally mean, and how far back does \emph{amāre} go? (One who does it \textbf{for love}; the trail \textbf{stops at Latin}.) Say the chapter's four verbs with \emph{je}. (\emph{Je prends, je demande, j'aide, j'aime}.)
\end{tcolorbox}
